% introduction.tex -- Introduction
% Three-paragraph structure: (1) problem, (2) landscape, (3) contribution.

Attention is the central computational primitive of modern deep learning~\cite{vaswaniAttentionAllYou2017}. Given a query, it computes a softmax-weighted average of stored values, a powerful operation but a fundamentally \emph{deterministic} one. The same query always produces the same output. Attention retrieves; it does not generate. Yet generation from a structured memory is precisely what many downstream tasks require: producing novel but plausible continuations, interpolating between stored prototypes, or exploring the space of patterns consistent with partial evidence. This limitation is sharpest in the low-data regime---small protein families, rare event catalogs, specialized image collections---where the memory bank contains tens to hundreds of examples, far too few to train a score network or diffusion model, yet enough to define a structured energy landscape. A natural question arises: \emph{can the attention mechanism itself be made stochastic, in a principled way, so that it samples from the space of memories rather than merely returning their weighted average?}

The ingredients for an answer already exist, though they have not been assembled. On one side, modern Hopfield networks~\cite{ramsauerHopfieldNetworksAll2021}, building on classical associative memory~\cite{hopfieldNeuralNetworksPhysical1982} and dense associative memories~\cite{krotovDenseAssociativeMemory2016,krotovLargeAssociativeMemory2021}, have revealed that each attention head implicitly performs gradient descent on a smooth, confining energy whose minima are stored patterns. On the other side, Langevin dynamics~\cite{robertsTweedieExponentialConvergence1996,wellingBayesianLearningStochastic2011,durmusNonasymptoticAnalysisUnadjusted2017} converts any such energy into a sampler for the corresponding Boltzmann distribution by adding calibrated noise to the gradient update. Energy-based models~\cite{lecunTutorialEnergyBased2006,songHowTrainEnergyBased2021} and score-based diffusion methods~\cite{songGenerativeModelingEstimating2019,songScoreBasedGenerativeModeling2021} have demonstrated the power of Langevin-type sampling for generation, but rely on black-box neural networks whose scores must be learned. The Energy Transformer~\cite{hooverEnergyTransformer2024} brings Hopfield energies into a deep architecture, yet remains a discriminative model that descends the energy rather than sampling from it. The classical retrieval-generation duality---Hopfield networks retrieve, Boltzmann machines sample, both over the same energy~\cite{ackeleyLearningAlgorithmBoltzmann1985}---has not been lifted to the modern continuous setting.

In this paper, we close that gap. We show that applying the unadjusted Langevin algorithm to the modern Hopfield energy yields a \emph{stochastic attention} update: a contraction toward the origin, a softmax attention pull toward stored memories, and an isotropic Gaussian perturbation governed by a temperature parameter. The inverse temperature $\beta$ interpolates between exact retrieval ($\beta\to\infty$) and open-ended generation ($\beta\to 0$), requiring no learned score network, no training loop, and no contrastive objective. Because the Hopfield energy gradient is exactly the identity minus the attention map, every step costs the same as a single attention head; in particular, $\mathbf{X}$ can be the key matrix of any pretrained attention layer, making stochastic attention a zero-shot stochastic decoding layer that requires no retraining. The energy's analytic structure (smoothness, Lipschitz gradients, quadratic confinement) delivers convergence guarantees unavailable to generic energy-based models. The contribution lies in the synthesis and the analysis it enables. First, the closed-form score function eliminates the score matching or contrastive divergence required by standard energy-based and score-based generative models. Second, by analyzing the Hopfield energy \emph{as a sampling target} rather than for retrieval alone, we derive an entropy inflection condition (Proposition~\ref{prop:entropy-derivative}) that identifies the critical $\beta^*$ at which the sampler transitions from retrieval to generation, with a scaling law $\beta^*\sim\sqrt{d}$ yielding a dimension-dependent SNR criterion for temperature selection. Third, we characterize the smoothness, dissipativity, and condition number of the energy as a function of $\beta$ and $K/d$ for sampling guarantees (Proposition~\ref{prop:regularity}). In practice, the same algorithm realizes two qualitatively distinct operating regimes -- \emph{structured retrieval} at high $\beta$ and \emph{genuine generation} at low $\beta$ -- producing, e.g., $6.9{\times}$ lower amino acid composition divergence than a VAE on protein sequences at matched novelty, despite using zero training. Section~2 surveys related work; Section~3 reviews the background; Section~4 derives the algorithm and its properties; Section~5 validates both regimes across five domains.
